Let \(c_i=1/e_i\) for \(i=1,2\).  If a single generator \(f\) were universally exact at both propensities, \(\mathrm{thm:one\text{-}sided\text{-}generator\text{-}iff\text{-}frontier}\) would give
\[
 f\in\mathfrak H_{c_1}\cap\mathfrak H_{c_2}. \tag{20}
\]
Assume without loss of generality that \(c_1<c_2\).  Membership in \(\mathfrak H_{c_i}\) supplies \(b_i\in\mathbb R\) such that
\[
 f(t)=b_i(t-1),\qquad 0\le t\le c_i.
\]
Because \([0,c_1]\) contains points other than \(1\), these two identities force \(b_1=b_2=:b\).  The representative \(g(t)=f(t)-b(t-1)\) therefore vanishes on \([0,c_2]\).  But membership in \(\mathfrak H_{c_1}\) requires \(g(t)>0\) for every \(t>c_1\), contradicting its vanishing on \((c_1,c_2]\).  Thus (20) is impossible.

For the positive result, fix \(e\in(0,1)\) and let \(f_e(t)=(t-1/e)_+\).  This function is finite, continuous, and convex, satisfies \(f_e(1)=0\), vanishes on \([0,1/e]\), and is strictly positive above \(1/e\).  Taking \(b=0\) in the definition gives \(f_e\in\mathfrak H_{1/e}\).  The one-sided frontier theorem therefore gives exactness at that stratum, and the mutual-absolute-continuity frontier theorem gives exactness in the mutual regime as well.